\begin{figure}[ht]
\centering
\begin{tikzpicture}[>=Stealth, font=\small, node distance=1.8cm]
  \node[draw, rounded corners, fill=blue!8, minimum width=4.6cm, minimum height=1cm] (tp)
    {$T_p$ 在 $A_f$ 上};
  \node[draw, rounded corners, fill=green!8, below left=of tp, minimum width=4.5cm, minimum height=1cm] (es)
    {$T_p=\Frob_p+p\Frob_p^{-1}$};
  \node[draw, rounded corners, fill=orange!10, below right=of tp, minimum width=4.8cm, minimum height=1cm] (eig)
    {$T_p$ 的特征值是 $a_p(f^\sigma)$};
  \node[draw, rounded corners, fill=purple!10, below=of $(es)!0.5!(eig)$, minimum width=5.2cm, minimum height=1cm] (poly)
    {$X^2-a_p(f^\sigma)X+p$};
  \node[draw, rounded corners, fill=red!8, below=of poly, minimum width=5.3cm, minimum height=1cm] (lp)
    {$L_p(s,A_f)=\prod_\sigma(1-a_p(f^\sigma)p^{-s}+p^{1-2s})^{-1}$};

  \draw[->, thick] (tp) -- (es);
  \draw[->, thick] (tp) -- (eig);
  \draw[->, thick] (es) -- (poly);
  \draw[->, thick] (eig) -- (poly);
  \draw[->, thick] (poly) -- (lp);
\end{tikzpicture}
\caption{Eichler--Shimura 关系把 Hecke 特征值与 Frobenius 特征多项式识别起来，因此 $A_f$ 的局部 $L$-因子由 $f$ 的 Fourier 系数直接给出。}
\label{fig:ch08-l-bridge}
\end{figure}
